\section{Canonical linear division in original flat Witt modules}
\label{sec:witt-two-module-division}

For the actual flat $W_2(k)$-module $M$, let $pM$
be the range of its original prime scalar map,
and let $\overline M=M/pM$ be the actual module quotient.

\begin{definition}[Original prime image and reduction]
\label{def:witt-two-module-prime-image}
\leanok
\lean{EnvelopeProperty.DeligneIllusie.WittTwoModuleDivision.primeSubmodule,EnvelopeProperty.DeligneIllusie.WittTwoModuleDivision.Reduction,EnvelopeProperty.DeligneIllusie.WittTwoModuleDivision.primeMultiple}
The original linear map $m\mapsto pm$ has actual
range submodule $pM$. Its value at each $m$,
regarded as an element of that range, is the prime
multiple of $m$. Reduction is the original quotient
module $M/pM$.
\uses{def:witt-two-module-prime}
\end{definition}
\begin{proof}
A linear map has a range submodule, and each of its
values belongs to that range with the original
argument as a preimage. The quotient by this
submodule is a module with its canonical linear
projection, giving all the stated constructions.
\end{proof}

\begin{theorem}[Canonical original linear division]
\label{thm:witt-two-module-linear-division}
\leanok
\lean{EnvelopeProperty.DeligneIllusie.WittTwoModuleDivision.dividePrimeEquiv,EnvelopeProperty.DeligneIllusie.WittTwoModuleDivision.dividePrime_primeMultiple}
There is an actual $W_2(k)$-linear equivalence
$\delta_M:pM\simeq\overline M$, whose value
on every original prime multiple $pm$ is the
original quotient class $\overline m$.
\uses{def:witt-two-module-prime-image,thm:witt-two-flat-prime-exact}
\end{theorem}
\begin{proof}
The first isomorphism theorem applied to the actual
map $\mu_M$ gives
$M/\ker\mu_M\simeq\operatorname{im}\mu_M=pM$,
sending $\overline m$ to $pm$. Prime exactness
identifies the actual submodules
$\ker\mu_M=pM$. The induced quotient identification
therefore gives an equivalence $M/pM\simeq pM$.
Take its inverse. On $pm$, the inverse first
isomorphism map returns the class of $m$ in
$M/\ker\mu_M$, and the quotient identification
returns exactly its original class in $M/pM$.
This proves the formula and linearity of the same map.
\end{proof}

\begin{definition}[The original maps on prime images and quotients]
\label{def:witt-two-module-prime-functoriality}
\leanok
\lean{EnvelopeProperty.DeligneIllusie.WittTwoModuleDivision.primeMap,EnvelopeProperty.DeligneIllusie.WittTwoModuleDivision.reductionMap}
Every actual $W_2(k)$-linear map $T:M\to N$
induces original linear maps $pT:pM\to pN$
and $\overline T:\overline M\to\overline N$.
\uses{def:witt-two-module-prime-image}
\end{definition}
\begin{proof}
For any $pm$, linearity gives $T(pm)=pT(m)$,
so $T$ preserves the actual prime image submodules.
Restricting $T$ therefore defines $pT$, and the
same submodule preservation induces the actual
quotient map $\overline T$.
\end{proof}

\begin{theorem}[Naturality of the same divided maps]
\label{thm:witt-two-module-division-naturality}
\leanok
\lean{EnvelopeProperty.DeligneIllusie.WittTwoModuleDivision.dividePrime_naturality}
If $M$ and $N$ are actual flat modules, every
original linear $T:M\to N$ satisfies
$\delta_N\circ pT=\overline T\circ\delta_M$
on the entire original prime image.
\uses{thm:witt-two-module-linear-division,def:witt-two-module-prime-functoriality}
\end{theorem}
\begin{proof}
Every element of $pM$ is $pm$ for an original
$m$. The left side at $pm$ is
$\delta_N(pT(m))=\overline{T(m)}$ by the original
value formula. The right side is
$\overline T(\overline m)=\overline{T(m)}$ by
the original induced quotient map. Thus the two
linear maps agree on every element of their domain.
\end{proof}
